\RequirePackage{pdfmanagement-testphase}
\DeclareDocumentMetadata{pdfversion=1.7,lang=en-US}
\documentclass[11pt]{amsart}

\usepackage[T1]{fontenc}
\usepackage{lmodern}
\usepackage{microtype}
\usepackage{amssymb}
\usepackage{booktabs}
\usepackage{longtable}
\usepackage{array}
\usepackage{hyperref}
\usepackage[nameinlink,capitalize]{cleveref}

\hypersetup{
  colorlinks=true,
  linkcolor=blue,
  citecolor=blue,
  urlcolor=blue,
  pdftitle={Nonabelian composition factors of finite groups whose maximal subgroups are complemented},
  pdfauthor={Richie Sater},
  pdfsubject={Conditional classification for Kourovka Problem 18.68},
  pdfkeywords={finite groups, composition factors, complemented maximal subgroups,
    primitive groups, maximal factorizations, Kourovka Notebook}
}

\newcommand{\CMP}{\mathsf{CMP}}
\newcommand{\Soc}{\operatorname{Soc}}
\newcommand{\Aut}{\operatorname{Aut}}
\newcommand{\Out}{\operatorname{Out}}
\newcommand{\Core}{\operatorname{Core}}
\newcommand{\Sym}{\operatorname{Sym}}
\newcommand{\PSL}{\operatorname{PSL}}
\newcommand{\PSU}{\operatorname{PSU}}
\newcommand{\PSp}{\operatorname{PSp}}
\newcommand{\Sp}{\operatorname{Sp}}

\newtheorem{theorem}{Theorem}[section]
\newtheorem{proposition}[theorem]{Proposition}
\newtheorem{lemma}[theorem]{Lemma}
\newtheorem{corollary}[theorem]{Corollary}
\theoremstyle{definition}
\newtheorem{definition}[theorem]{Definition}
\newtheorem{remark}[theorem]{Remark}

\title[Complemented maximal subgroups]{Nonabelian composition factors of finite groups\\
whose maximal subgroups are complemented}
\author{Richie Sater}
\date{Preprint, 31 August 2026}
\keywords{Finite groups, composition factors, complemented maximal subgroups,
primitive groups, maximal factorizations, Kourovka Notebook}

\begin{document}
\raggedbottom

\begin{abstract}
Assume the classification of finite simple groups and the external
factorization consequences recorded in Appendix~\ref{app:dependencies}.
Let \(G\) be a finite group in which every maximal subgroup \(M\) has an
exact complement: \(G=MK\) and \(M\cap K=1\) for some \(K\leq G\).  We prove
that every nonabelian composition factor of \(G\) is isomorphic to
\[
\PSL_2(7),\qquad \PSL_2(11),\qquad\text{or}\qquad \PSL_5(2),
\]
and that all three possibilities occur.  Subject to those inputs, this
answers Kourovka Problem 18.68.  Two load-bearing LPS source matches remain
at the explicit verification boundary described in the introduction and
Appendix~\ref{app:dependencies}.

The central difficulty is that the complement property does not pass to
sections.  We replace that invalid inheritance step by a section-safe
construction: a chosen composition factor \(S\) yields a monolithic quotient
with socle \(S^k\), and its almost-simple coordinate action lifts to a
faithful primitive action on a Cartesian power.  A regular complement would
then force a core-free transitive subgroup in the coordinate action.
Factor-free, prime-elusive, and socle-valuation obstructions reduce this
coordinate obligation to published exhaustive factorization
classifications.  The only infinite screen survivor,
\(\PSp_4(2^f)\), is eliminated by a prime-degree subfield action and a
2-adic argument.  Complete tables of marks certify only fifteen named small
almost-simple cases; they are not used for any infinite-family conclusion.
\end{abstract}

\maketitle

\section{Introduction}\label{sec:introduction}

Say that a finite group \(G\) has property \(\CMP\) if every maximal subgroup
\(M<G\) has a subgroup \(K\leq G\) such that
\begin{equation}\label{eq:complement}
G=MK,\qquad M\cap K=1.
\end{equation}
The complement is not assumed normal.  Kourovka Problem 18.68 asks for the
nonabelian composition factors of a nonsoluble group with this property
\cite[Problem 18.68]{kourovka21}.

Levchuk and Likharev proved that the nonabelian simple groups having
\(\CMP\) are precisely \(\PSL_2(7),\PSL_2(11),\PSL_5(2)\)
\cite[Theorem 1 and Lemmas 4, 6]{levchuk-likharev}.  Maslova obtained the same
composition-factor list under the stronger hypothesis that every maximal
subgroup is a Hall subgroup \cite[Theorem 1]{maslova-hall-factors}.  Maslova
and Revin subsequently described groups under that hypothesis and proved
that it implies \(\CMP\) \cite[Theorems 1--2]{maslova-revin}; the implication
is not reversible, so those results do not answer the present question.  The
complete classification of exact factorizations of almost-simple groups
\cite{li-wang-xia} likewise addresses the one-coordinate setting, not an
ambient monolithic group with socle \(S^k\).  Recent work of
Monakhov--Sokhor concerns monotone index sequences in maximal chains
\cite{monakhov-sokhor}, rather than exact complements or the nonabelian
composition factors considered here.  A literature search updated on
31 August 2026 found no published solution of Problem 18.68; this negative
search is not a claim of absolute priority.  The July 2026 edition of the
Kourovka Notebook still prints the problem without a solution annotation.

The structural setting has established roots.  Aschbacher--Scott organize
maximal subgroups through minimal normal subgroups
\cite{aschbacher-scott}; Kov\'acs studies maximal subgroups of composite
finite groups \cite{kovacs-composite} and primitive subgroups of wreath
products in product action \cite{kovacs-product}.  We do not claim the
chief-factor or product-action frameworks themselves.  The contribution here
is their CMP-specific synthesis: quotient closure places an arbitrary
composition factor in a monolithic group, a selected coordinate action lifts
to that same quotient, and three complementary obstructions close the
resulting classification screen.

The obstacle is inheritance.  Property \(\CMP\) passes to quotients but not
to normal subgroups, hence not automatically to sections.  For example,
\(D_8\) has \(\CMP\), while its normal cyclic subgroup \(C_4\) does not.
Thus one cannot attach \(\CMP\) directly to a nonabelian composition factor
and quote the simple-group theorem.

Our CMP-specific replacement is the section-safe product-action lift in
\cref{thm:product-lifting}.  Starting from a chosen composition factor \(S\),
we pass only to a quotient \(L\) whose unique minimal normal subgroup is
\(S^k\).  The group induced by a coordinate is almost simple,
\(S\leq X\leq\Aut(S)\).  A suitable primitive action
\(X\curvearrowright\Delta\) then manufactures a faithful primitive action
\(L\curvearrowright\Delta^k\).  If \(L\) retained \(\CMP\), its point
stabilizer would have a regular complement.  Under the first daggered
consequence attributed to Liebeck--Praeger--Saxl (LPS), that regular subgroup
forces a coordinate factorization, where three obstruction modes can be
applied.

\begin{theorem}\label{thm:main}
Assume CFSG and the external classification and factorization statements
listed in \cref{app:dependencies}, including the two LPS consequences marked
there with a dagger.
Let \(G\) be a finite group with \(\CMP\).  Every nonabelian composition
factor of \(G\) is isomorphic to one of
\[
\PSL_2(7),\qquad \PSL_2(11),\qquad \PSL_5(2).
\]
Each of the three possibilities occurs.
\end{theorem}

The two daggered inputs are LPS (2000), Corollary~3(iv), for the
product-action coordinate consequence, and LPS (1990), Theorem~D with
Remark~2, for the large alternating branch.  Their exact primary texts were
not independently reopened during the source traces through 31 August 2026.
This paper therefore consumes only the precise consequences stated in
\cref{app:dependencies}; if either source has narrower hypotheses, the
corresponding branch is unchecked.

The reusable CMP-specific part of the proof is the lift from a quotient-closed
maximal-subgroup property and the three-mode coordinate toolkit, not the
underlying product-action framework or the subsequent enumeration of families
supplied by the classification of finite simple groups (CFSG).  We therefore
prove the structural reduction in
\cref{sec:lifting} and \cref{sec:obstructions}, give one instructive use of each
obstruction in \cref{sec:representative}, and state the exhaustive
classification interface in \cref{sec:classification}.  The main theorem is
assembled in \cref{sec:main-proof}.  Repetitive family checks and finite
certificate details are placed in \cref{app:classification} and \cref{app:computer};
\cref{app:dependencies} records the exact external inputs and trust boundary.

\section{Proof architecture and the section-safe lift}\label{sec:lifting}

The logical route from a composition factor to a contradiction is
\begin{equation}\label{eq:architecture}
\boxed{\begin{gathered}
S\ \leadsto\ L\text{ with }\Soc(L)=S^k
\ \leadsto\ X=N_L(S_1)/C_L(S_1),\\
L\hookrightarrow X\wr P\curvearrowright\Delta^k
\ \leadsto\ L_\omega\text{ with no complement}.
\end{gathered}}
\end{equation}
Here \(\leadsto\) denotes a construction, not a homomorphism.  In particular,
\(X\) is the group induced by a coordinate stabilizer; the proof never
asserts a quotient map \(L\to X\).  This distinction is what makes the
argument safe when \(\CMP\) fails to pass to sections.

\subsection{The valid inheritance step}

\begin{lemma}[quotient closure]\label{lem:quotient}
If \(G\) has \(\CMP\) and \(N\trianglelefteq G\), then \(G/N\) has \(\CMP\).
\end{lemma}

\begin{proof}
Let \(\overline M\) be maximal in \(G/N\), and let \(M\) be its full
preimage.  Choose \(K\leq G\) satisfying \eqref{eq:complement}.  Then
\(G/N=\overline M(KN/N)\).  If \(kN\in\overline M\) with \(k\in K\), then
\(k\in M\), and hence \(k=1\).  Thus \(KN/N\) is an exact complement to
\(\overline M\).
\end{proof}

\begin{proposition}[monolithic reduction]\label{prop:monolithic}
Let \(S\) be a nonabelian composition factor of a finite \(\CMP\)-group
\(G\).  There is a quotient \(L\) of \(G\) such that
\begin{enumerate}
\item \(L\) has \(\CMP\);
\item \(\Soc(L)=N\cong S^k\) for some \(k\geq1\), and \(N\) is the unique
minimal normal subgroup of \(L\);
\item \(C_L(N)=1\).
\end{enumerate}
\end{proposition}

\begin{proof}
Choose a chief factor \(A/B\cong S^k\) of \(G\), and define
\[
C/B=C_{G/B}(A/B),\qquad L=G/C,\qquad N=AC/C.
\]
Since \(A/B\) is centerless, \(A\cap C=B\), and hence \(N\cong A/B\).
Chiefness makes \(N\) a minimal normal subgroup of \(L\).  If \(gC\)
centralizes \(AC/C\), then \([g,A]\leq A\cap C=B\), so \(g\in C\) by the
definition of \(C/B=C_{G/B}(A/B)\).  Thus the conjugation action of \(L\) on
\(N\) is faithful and \(C_L(N)=1\).  If \(R\) were another minimal normal subgroup, then
\([R,N]\leq R\cap N=1\), forcing \(R\leq C_L(N)=1\).  Thus \(N\) is unique.
Finally, \(L\) has \(\CMP\) by \cref{lem:quotient}.
\end{proof}


The example \(C_4\trianglelefteq D_8\) shows why the preceding proposition,
rather than section inheritance, is the correct starting point.  The next
construction turns its entire direct-power socle into a coordinate problem.

\subsection{Manufacturing the product action}

Write
\[
N=S_1\times\cdots\times S_k\cong S^k.
\]
The action of \(L\) on the factors is transitive, since the product of the
factors in any orbit is normal in \(L\).  Set
\begin{equation}\label{eq:coordinate}
X=N_L(S_1)/C_L(S_1),
\end{equation}
identifying \(S_1\) with its inner automorphism group.  Then
\(S\leq X\leq\Aut(S)\).  The standard wreath embedding gives
\begin{equation}\label{eq:wreath}
L\hookrightarrow X\wr P\leq X\wr\Sym(k),
\end{equation}
where \(P\) is transitive, \(N\) is the base socle \(S^k\), and the group
induced on a coordinate is \(X\).  An element in the kernel fixes every
factor and induces the identity automorphism on each \(S_i\), so it
centralizes \(N\).  Thus the kernel lies in \(C_L(N)=1\), and the embedding
is faithful.

Let \(H<X\) be core-free and maximal.  Put
\(\Delta=X/H\) and \(V=H\cap S\), and assume \(V\ne1\).  Since \(S\) is
normal and is not contained in \(H\), we have \(X=SH\); in particular,
\(S\) is transitive on \(\Delta\).  Via \eqref{eq:wreath}, let \(L\) act on
\(\Delta^k\) in product action.

\begin{theorem}[section-safe product-action lifting]\label{thm:product-lifting}\label{lem:pa-primitive}
The action of \(L\) on \(\Delta^k\) is faithful and primitive.  If
\(\omega=(H,\ldots,H)\) and \(M=L_\omega\), then
\(L=MS^k\) and \(M\cap S^k=V^k\).
\end{theorem}

Here we apply standard wreath-product and Scott-lemma ingredients.  The proof
tests an intermediate overgroup through its intersection with \(S^k\), then
uses Scott's lemma to show that the only possible strips have length one.

\begin{proof}
Let \(\omega=(H,\ldots,H)\), \(M=L_\omega\), and \(N=S^k\).  The group \(N\)
is transitive, so \(L=MN\) and \(M\cap N=V^k\).

Suppose \(M\leq J\leq L\), and put \(D=J\cap N\).  Then \(J=MD\),
\(D\trianglelefteq M\), and \(V^k\leq D\).  The image of \(M\) on the
coordinates is transitive: multiply an element of \(L\) by a suitable base
element to make it fix \(\omega\) without changing its top permutation.
Moreover, the subgroup of \(M\) stabilizing coordinate \(i\) induces \(H\)
there.  For a prescribed \(h\in H\), lift \(h\) through
\(N_L(S_i)\to X\).  Its \(i\)-coordinate already fixes the coset \(H\), and
the transitivity of each other factor permits multiplication by suitable
elements of \(S_j\), one for each \(j\ne i\), that return all remaining
coordinates to \(H\).  The corrected lift lies in \(M\) and still induces
\(h\) on \(S_i\).

Let \(D_i\) be the projection of \(D\) to \(S_i\cong S\).  It contains \(V\)
and is normalized by \(H\).  Maximality of \(H\) implies either \(D_i=V\) or
\(HD_i=X\).  In the second case
\[
D_i=(H\cap S)D_i=(HD_i)\cap S=X\cap S=S.
\]
Coordinate transitivity gives the same alternative for every \(i\).  If all
projections are \(V\), then \(D=V^k\) and \(J=M\).  Otherwise \(D\) is
subdirect in \(S^k\).  By Scott's lemma
\cite[Lemma, p.~328]{scott} it is a direct product of full
diagonal strips.  Since \(V^k\leq D\) and \(V\ne1\), every strip has length
one, so \(D=N\) and \(J=L\).  Thus \(M\) is maximal.  Faithfulness follows
from core-freeness of \(H\) and \eqref{eq:wreath}.
\end{proof}


The theorem is the decisive removal of the unbounded exponent \(k\).  It
constructs a new action of the same quotient \(L\); it does not assume that
the original primitive action supplied by \(\CMP\) already has product-action
type.  The next section explains how the existence of a complement in this
manufactured action constrains the one-coordinate group \(X\).

\section{The coordinate-obstruction toolkit}\label{sec:obstructions}

Let \(H<X\) be as in \cref{thm:product-lifting}, put
\(\Delta=X/H\), and write \(V=H\cap S\ne1\).  A complement to the lifted
point stabilizer is a regular subgroup, so it can be ruled out at any of
three increasingly flexible levels.

Throughout, \(X/H\) denotes the set \(\{Hx:x\in X\}\) of right cosets, on
which \(X\) acts by right multiplication.  Dotted Atlas notation such as
\(M_{12}.2\) denotes an extension without asserting that it splits; a colon
denotes a split extension, and subscripts distinguish Atlas isomorphism
types.

\begin{center}
\footnotesize
\begin{tabular}{>{\raggedright\arraybackslash}p{.20\textwidth}
                >{\raggedright\arraybackslash}p{.35\textwidth}
                >{\raggedright\arraybackslash}p{.35\textwidth}}
\toprule
Mode & Coordinate statement & Lifting mechanism \\
\midrule
factor-free & no core-free \(C<X\) satisfies \(X=HC\)
& the daggered LPS Corollary 3(iv) consequence, with a direct almost-simple
argument for \(k=1\) \\
prime-elusive & \(p\mid\lvert\Delta\rvert\) and every element of order \(p\)
fixes a point & fixed points lift to the full wreath product on \(\Delta^k\) \\
socle valuation & \(S\) is \(p\)-elusive and
\(v_p(\lvert\Delta\rvert)>v_p(\lvert X:S\rvert)\)
& a regular subgroup has a nontrivial \(p\)-element in its intersection with
\(S^k\) \\
\bottomrule
\end{tabular}
\end{center}

The formal statements below isolate exactly what the classification appendix
must supply.

\begin{lemma}[regular-action reformulation]\label{lem:regular}
A subgroup \(K\leq G\) satisfies \eqref{eq:complement} if and only if it is
regular on the right cosets of \(M\).  If \(M\) is maximal and core-free,
this is a faithful primitive action.
\end{lemma}

\begin{proof}
Transitivity is equivalent to \(G=MK\), and the stabilizer in \(K\) of the
coset \(M\) is \(M\cap K\).
\end{proof}

Call \((X,\Delta)\) \emph{factor-free} if it has no core-free transitive
subgroup; equivalently, there is no core-free \(C<X\) with \(X=HC\).

\begin{proposition}[coordinate obstruction]\label{prop:coordinate}
If \((X,X/H)\) is factor-free, then \(L\) does not have \(\CMP\).
\end{proposition}

\begin{proof}
For \(k=1\), we have \(L=X\).  If a regular subgroup \(R\) existed, then
\(R\) would be transitive and \(R\cap H=1\).  Its core in \(X\), if
nontrivial, would contain the socle \(S\), as \(X\) is almost simple.  But
then \(S\leq R\), contrary to
\(1\ne H\cap S\leq H\) and \(R\cap H=1\).  Thus \(R\) would be core-free,
contradicting factor-freeness.  Let \(k\ge2\).  The action in
\cref{thm:product-lifting} is a
primitive product action with almost-simple component
\((X,X/H)\): indeed, its socle stabilizer is \(V^k\), which is nontrivial
and is not subdirect in \(S^k\).  The first daggered source assumption is
the consequence attributed to Corollary 3(iv) of \cite{lps-transitive}: if
such a group has a regular subgroup, then its component action has a core-free
transitive subgroup.  Hence no regular subgroup exists.  But if \(L\) had
\(\CMP\), its maximal point stabilizer would have a complement, regular by
\cref{lem:regular}, a contradiction.
\end{proof}

This is the step that removes the unbounded exponent \(k\): it converts the
entire product-action problem into one almost-simple factorization screen.

We first state the mechanism for applying maximal-factorization tables.

\begin{lemma}[factor screen]\label{lem:factor-screen}
Let \(1<V<S\), and assume that the \(S\)-conjugacy class of \(V\) is
invariant under \(X\).  Choose a maximal subgroup \(H<X\) containing
\(N_X(V)\).  Then \(H\) is core-free and \(X=HS\).  If \(X/H\) is not
factor-free, there are \(S\leq Y\leq X\) and a nontrivial maximal
factorization \(Y=AB\) such that, after interchanging the factors,
\[
V\leq A\cap S.
\]
\end{lemma}

\begin{proof}
Class invariance gives \(N_X(V)S=X\), so no proper overgroup of \(N_X(V)\)
contains \(S\).  Hence \(H\) is core-free and \(HS=X\).

Suppose \(C<X\) is core-free and \(X=HC\).  Put \(Y=CS\).  Then
\(Y=(H\cap Y)C\), and both factors supplement \(S\) in \(Y\).  Consequently
all their proper maximal overgroups are core-free.  Enlarging them gives a
maximal factorization \(Y=AB\), with
\(V\leq H\cap S\leq A\cap S\).
\end{proof}

\begin{definition}
An action of \(Z\) on \(\Lambda\) is \emph{\(p\)-elusive} if
\(p\mid\lvert\Lambda\rvert\) and every element of order \(p\) in \(Z\)
fixes a point.
\end{definition}

\begin{lemma}[wreath lifting]\label{lem:elusive-wreath}
If \(X\) is \(p\)-elusive on \(\Delta\), then
\(X\wr\Sym(k)\) is \(p\)-elusive on \(\Delta^k\).  In particular, no
subgroup of this wreath product is regular on \(\Delta^k\).
\end{lemma}

\begin{proof}
Let \(w=(x_1,\ldots,x_k)\sigma\) have order \(p\).  On a fixed coordinate of
\(\sigma\), the corresponding \(x_i\) has order \(1\) or \(p\), and hence
fixes a point.  On a \(p\)-cycle of \(\sigma\), the cycle product of the
\(x_i\) is \(1\), since \(w^p=1\); the fixed-tuple equations can therefore
be solved successively around the cycle.  Thus \(w\) fixes a tuple.

If a subgroup were regular, Cauchy's theorem would give it a nonidentity
element of order \(p\), because its order is
\(\lvert\Delta\rvert^k\).  That element would fix a point, a contradiction.
\end{proof}

\begin{lemma}[socle-valuation lifting]\label{lem:valuation}
Suppose \(S\) is \(p\)-elusive on \(\Delta=X/H\).  Let
\[
a=v_p(\lvert\Delta\rvert),\qquad
o=v_p(\lvert X:S\rvert).
\]
If \(a>o\), then every subgroup \(L\leq X\wr\Sym(k)\) containing \(S^k\)
contains no regular subgroup on \(\Delta^k\).
\end{lemma}

\begin{proof}
Suppose \(R\leq L\) is regular and put \(Q=R\cap S^k\).  Since
\(R/Q\hookrightarrow L/S^k\),
\[
v_p(\lvert R:Q\rvert)\leq ko+v_p(k!).
\]
As \(v_p(\lvert R\rvert)=ka\),
\[
v_p(\lvert Q\rvert)\geq k(a-o)-v_p(k!)>0,
\]
because \(a-o\ge1\) and \(v_p(k!)<k\).  An element of order \(p\) in \(Q\)
fixes a tuple coordinatewise, contrary to regularity.
\end{proof}


\section{Representative obstruction proofs}\label{sec:representative}

This section gives one complete example of each mode.  The point is to show
how the toolkit works; the repetitive family coverage is deferred to
\cref{app:classification}.

\subsection{\texorpdfstring{Prime-elusive: two-subsets for \(A_5\)}%
{Prime-elusive: two-subsets for A5}}

\begin{proposition}\label{prop:a5-elusive}
For \(A_5\leq X\leq S_5\), the action of \(X\) on the ten 2-subsets of
\(\{1,\ldots,5\}\) is \(2\)-elusive.  Its point stabilizer is core-free and
maximal, and has nontrivial intersection with \(A_5\).
\end{proposition}

\begin{proof}
Write \(H\) for the stabilizer of a 2-subset.  It is maximal in \(S_5\),
and \(H\cap A_5\) is maximal in \(A_5\); the actions are therefore
primitive and faithful.
Every involution in \(S_5\), and hence in either coordinate group, contains a
2-cycle and fixes its support as a 2-subset.  Since \(2\mid10\), the action is
2-elusive.  The subgroup \(H\cap A_5\) contains a double transposition and is
nontrivial.
\end{proof}

By \cref{lem:elusive-wreath}, this single fixed-point observation excludes a
regular complement on every Cartesian power.

\subsection{Factor-free: six-subsets for large alternating groups}

\begin{proposition}\label{prop:alternating-screen}
Let \(n\ge13\), let \(A_n\leq X\leq S_n\), and let \(H\) stabilize a
6-subset.  Then \(H\) is core-free and maximal, \(H\cap A_n\ne1\), and
\(X/H\) is factor-free.
\end{proposition}

\begin{proof}
The stabilizer is maximal because \(6<n/2\).  The second daggered source
assumption, attributed to the alternating-group part of the LPS
maximal-factorization theorem
\cite[Theorem D and Remark 2]{lps-maxfact}, is that an intransitive factor
in a nontrivial maximal factorization stabilizes a \(j\)-subset with
\(j\le5\), apart from listed cases of degrees \(6,8,10\).  Consequently no
core-free transitive subgroup can factor \(X\) with the 6-set stabilizer.
Core-freeness follows because the only nontrivial normal subgroup of either
coordinate group contains \(A_n\), whereas the stabilizer does not; its
intersection with \(A_n\) contains \(A_6\times A_{n-6}\).
\end{proof}

The factor screen converts the published list into a containment test: the
chosen subgroup type never appears on either side of an allowed maximal
factorization.

\subsection{Socle valuation: the infinite symplectic survivor}

\begin{proposition}\label{prop:sp4-screen}
Let \(S=\PSp_4(2^f)\), \(f\ge2\), and
\(S\leq X\leq\Aut(S)\).  There is a core-free maximal subgroup \(H<X\)
with \(H\cap S\ne1\) such that \(S\) is 2-elusive on \(X/H\) and
\[
v_2(\lvert X:H\rvert)>v_2(\lvert X:S\rvert).
\]
Consequently the lifted product action has no regular subgroup for any
socle multiplicity \(k\ge1\).
\end{proposition}

\begin{proof}
Let \(q=2^f\).  Choose a prime \(r\mid f\), put
\(d=f/r\), \(q_0=2^d\), and let
\[
V=\Sp_4(q_0)<S
\]
be the standard subfield subgroup.  It is maximal in \(S\), since the
field-extension degree \(r\) is prime
\cite[Table 8.14]{bhrd}; see also
\cite[Proposition 4.2 and Table 3]{burness-base}.

We check explicitly that this one choice works for every field-only and
graph-outer coordinate group \(S\leq X\leq\Aut(S)\).  Realize \(S\) as the
fixed points of the standard Frobenius \(\varphi^f\) in the adjoint algebraic
group of type \(B_2\) over \(\overline{\mathbb F}_2\).  The exceptional
graph-field endomorphism \(\rho\) satisfies \(\rho^2=\varphi\), as in the
setup preceding Lemma~2.1 and Theorem~4(a) of \cite{harper}, while
\(\Aut(S)=\langle S,\rho\rangle\) by \cite[Lemma 2.1]{harper}.  Since
\(\rho\) commutes with \(\varphi\), it
normalizes the fixed-point subgroup
\(V=\operatorname{Sp}_4(2^d)\).  Thus the \(S\)-class of \(V\) is invariant
under all of \(\Aut(S)\).

Put \(H=N_X(V)\).  Class invariance gives \(X=SH\), and
\(H\cap S=N_S(V)=V\).  Moreover, \(H\) is maximal in \(X\).  Indeed, if
\(H<K<X\), then maximality of \(V\) in \(S\) gives
\(K\cap S=V\) or \(S\).  In the first case, \(V=K\cap S\trianglelefteq K\),
so \(K\leq N_X(V)=H\); in the second, \(S\leq K\) and
\(X=SH\leq K\).  Both are contradictions.  Finally, \(H\) is core-free:
every nontrivial normal subgroup of the almost-simple group \(X\) contains
\(S\), whereas \(H\cap S=V<S\).

Every involution class of \(S\) meets \(V\).  Here are the details needed
at the prime-field endpoint \(q_0=2\).  By Lemmas 2.1--2.2 and Corollary 2.3
of \cite{lps-regular}, the three involution classes of \(S\) meet
\[
 O_4^+(q)\cong
 (\operatorname{SL}_2(q)\times\operatorname{SL}_2(q))
 \rtimes\langle\tau\rangle,
\]
where \(\tau\) interchanges the factors, and conjugacy inside this orthogonal
group agrees with conjugacy in \(S\).  If \(u\) represents the unique
nonidentity involution class in \(\operatorname{SL}_2(q)\), the three
orthogonal classes have representatives
\((u,1),(u,u),\tau\).  We may take
\(u\in\operatorname{SL}_2(2)\), so all three representatives lie in
\(O_4^+(2)\leq\operatorname{Sp}_4(2)\leq V\).  Hence \(S\) is
2-elusive on \(X/H\).

Since \(X=SH\), the degree and its 2-part are
\begin{equation}\label{eq:sp-degree}
 \lvert X:H\rvert=\lvert S:V\rvert
 =\frac{q^4(q^2-1)(q^4-1)}
 {q_0^4(q_0^2-1)(q_0^4-1)},\qquad
 a=v_2(\lvert X:H\rvert)=4(f-d).
\end{equation}
The same outer-automorphism description gives
\(\lvert X:S\rvert\mid2f\), so
\[
 o=v_2(\lvert X:S\rvert)\leq1+v_2(f)<2f\leq4(f-d)=a.
\]
Here \(v_2(f)\leq f-1\), and \(f-d=f(1-1/r)\geq f/2\).
Thus \cref{lem:valuation} applies uniformly.  In particular, when
\(q=4\) and \(X=\Aut(S)\), the corrected subgroup has
\(\lvert H\rvert=2880\), \(\lvert H\cap S\rvert=720\), and degree \(1360\); these values are
also reproduced in \cref{app:computer}.
\end{proof}

The valuation proof illustrates why a screen survivor is not a gap: the
weaker socle-level fixed-point statement becomes decisive once the outer
\(2\)-part is smaller than the degree \(2\)-part.  The Appendix A audit shows
that no other infinite family survives the first two modes.

\section{How the classification closes}\label{sec:classification}

\begin{proposition}\label{prop:classification}
Let \(S\) be a nonabelian finite simple group not isomorphic to
\(\PSL_2(7),\PSL_2(11)\), or \(\PSL_5(2)\), and let
\(S\leq X\leq\Aut(S)\).  There is a core-free maximal subgroup \(H<X\),
with \(H\cap S\ne1\), for which one of the following holds:
\begin{enumerate}
\item \(X/H\) is factor-free;
\item \(X\) is \(p\)-elusive on \(X/H\) for some prime \(p\);
\item \(S=\PSp_4(2^f)\), \(f\ge2\), and the hypotheses of
\cref{lem:valuation} hold with \(p=2\).
\end{enumerate}
\end{proposition}

\begin{proof}
The common proof template is to choose an automorphism-invariant subgroup
\(1<V<S\), place \(N_X(V)\) in a core-free maximal subgroup \(H\), and apply
\cref{lem:factor-screen}.  The following roadmap identifies every family,
its exceptional treatment, and the exact appendix location where the
containment obligation is checked.

\begin{center}
\small
\begin{tabular}{>{\raggedright\arraybackslash}p{.20\textwidth}
                >{\raggedright\arraybackslash}p{.37\textwidth}
                >{\raggedright\arraybackslash}p{.31\textwidth}}
\toprule
CFSG family & common screen and exceptions & complete audit \\
\midrule
alternating & small degrees by explicit elusive actions or five finite
certificates; \(n\ge13\) by 6-subsets &
\cref{prop:a5-elusive} and \cref{prop:alternating-screen}; \cref{app:alternating} \\
classical & invariant parabolics, torus normalizers, or geometric subgroups;
\(L_3(4)\) is finite and even \(\PSp_4\) uses valuation &
\cref{prop:sp4-screen}; \cref{app:classical} and \cref{app:computer} \\
exceptional Lie type & defining-characteristic Sylow or graph-stable
parabolic intersections & \cref{app:exceptional} \\
sporadic & Atlas maximal subgroups checked against Giudici's exhaustive
factorization tables; \(M_{24}\) is 2-elusive & \cref{app:sporadic} \\
\bottomrule
\end{tabular}
\end{center}

The alternating, classical, exceptional, and sporadic groups exhaust the
nonabelian finite simple groups.  The cited appendix rows include every outer
coordinate group and low-parameter isomorphism.  Together with the three
representative proofs, they establish the proposition.
\end{proof}

The classification is used only through this proposition.  Thus a first
reading may skip \cref{app:classification}: the structural reason for the
result is already visible, while every family obligation remains traceable.

\section{Proof of the main theorem}\label{sec:main-proof}

\begin{proof}[Proof of \cref{thm:main}]
Let \(S\) be a nonabelian composition factor of \(G\).  Apply
\cref{prop:monolithic} to obtain a \(\CMP\)-quotient \(L\) with unique
self-centralizing socle \(S^k\), and define its coordinate group \(X\) by
\eqref{eq:coordinate}.

Suppose \(S\) is not one of the three groups in the conclusion.  Choose the
maximal subgroup \(H<X\) supplied by \cref{prop:classification}.  By
\cref{thm:product-lifting}, the same group \(L\) has a faithful primitive action
on \((X/H)^k\).  If the action \(X/H\) is factor-free,
\cref{prop:coordinate} contradicts \(\CMP(L)\).  If it is \(p\)-elusive,
\cref{lem:elusive-wreath} shows that the primitive action has no regular
subgroup, again contradicting \(\CMP(L)\).  In the final symplectic case,
the same contradiction follows from \cref{lem:valuation}.  Therefore \(S\)
is one of the three displayed groups.

Conversely, Levchuk--Likharev prove that each of the three simple groups has
\(\CMP\) \cite[Lemmas 4 and 6]{levchuk-likharev}; taking the group itself shows
that each occurs as a composition factor.
\end{proof}


\section{Conclusion}\label{sec:yield}

The proof teaches how to constrain a composition factor when a maximal-
subgroup property is quotient-closed but not section-closed.  The monolithic
quotient keeps the ambient hypothesis intact, while
\cref{thm:product-lifting} converts its direct-power socle into a family of
coordinate actions.  This use of established monolithic and product-action
machinery is CMP-specific: it replaces an invalid inheritance assertion by a
primitive action whose point stabilizer is still a maximal subgroup of the
original quotient.

The essential hypotheses are visible in the proof.  Quotient closure is
needed to reach \(L\); the unique self-centralizing socle gives a faithful
wreath embedding; a nontrivial coordinate intersection \(H\cap S\) forces
primitivity; and the maximal-subgroup property must turn the lifted point
stabilizer into a regular subgroup.  Other quotient-closed properties that
force such a regular or suitably transitive partner may admit a similar
reduction, but no general metatheorem is claimed here.

The price of the final list is classification dependence.  The product-action
lift and the three obstruction lemmas do not themselves identify all simple
socles; the audited published lists, together with the two explicit LPS
source assumptions, supply that coverage.  A classification-free proof, or an axiomatization with a second
application, remains a natural direction.

\begin{center}
\small
\begin{tabular}{>{\raggedright\arraybackslash}p{.22\textwidth}
                >{\raggedright\arraybackslash}p{.66\textwidth}}
\toprule
Evidence layer & Exact role in \cref{thm:main} \\
\midrule
ordinary mathematics & quotient closure, monolithic reduction, the
product-action lift, all three obstruction lemmas, containment deductions,
the \(\PSp_4(2^f)\) subfield and valuation argument, and theorem assembly \\
published classifications & CFSG family exhaustiveness, the audited maximal-
factorization lists, and the stated subfield, automorphism, and
involution-class inputs \\
explicit source assumptions & the two precise consequences attributed to
LPS (2000), Corollary 3(iv), and LPS (1990), Theorem D with Remark 2 \\
GAP/TomLib & exact factor-free predicates for five small alternating
coordinate groups and ten groups with socle \(\PSL_3(4)\), plus nine
separately labelled sporadic cross-checks \\
Python integrity checks & pinned serialization, inventory, hashes, and
arithmetic; they do not repeat TomLib's intersection search \\
GAP/AtlasRep & a non-proof \(q=4\) regression for the symplectic
construction; it is not used to prove an infinite-family statement \\
formal proof assistants & none; the theorem is not claimed to have a
kernel-checked formal proof \\
\bottomrule
\end{tabular}
\end{center}

The layers compose in one direction: the ordinary reduction produces a
coordinate obligation; published classifications make the infinite screen
exhaustive; complete tables of marks settle only the fifteen named finite
residues.  \Cref{app:dependencies} gives theorem numbers and hypothesis
matches, and \cref{app:computer} gives the finite proposition checked by the
software.

\appendix
\section{The exhaustive coordinate-family audit}\label{app:classification}

This appendix proves the family assertions consumed by
\cref{prop:classification}.  Each subsection states the common screen,
identifies exceptional coordinate groups, and points to the exact published
classification locus.  The body retains the three arguments that teach the
method; this appendix retains the exhaustive case obligations.

\subsection{Alternating groups}\label{app:alternating}

For \(n\ne6\), the coordinate group is \(A_n\) or \(S_n\).  The following
actions are \(p\)-elusive:
\begin{center}
\begin{tabular}{lll}
\toprule
Socle & action and degree & prime \\
\midrule
\(A_5\) & 2-subsets, \(10\) & \(2\) \\
\(A_6\), \(X=A_6\) & natural, \(6\) & \(2\) \\
\(A_6\), \(X=S_6\) & unordered \(3+3\) partitions, \(10\) & \(2\) \\
\(A_8\) & 2-subsets, \(28\) & \(2\) \\
\(A_9\) & 2-subsets, \(36\) & \(2\) \\
\(A_{10}\) & 4-subsets, \(210\) & \(2\) \\
\(A_{11}\) & 3-subsets, \(165\) & \(3\) \\
\(A_{12}\) & 2-subsets, \(66\) & \(2\) \\
\bottomrule
\end{tabular}
\end{center}
These assertions follow from the cycle supports.  An involution fixes the
support of one 2-cycle.  In degree 10, every involution also has an invariant
4-set: use two 2-cycles, or one 2-cycle and two fixed points.  An element of
order 3 in degree 11 fixes the support of one 3-cycle.  For the \(S_6\)
partition action, the cycle types \(2,2^2,2^3\) respectively preserve a
3-set or interchange it with its complement.

The three outer extensions of \(A_6\) not equal to \(S_6\), and the two
groups \(A_7,S_7\), have factor-free maximal subgroups certified in
\cref{app:computer}.  Their orders and indices are
\[
\begin{array}{c|ccccc}
X&A_6.2_2&A_6.2_3&A_6.2^2&A_7&S_7\\ \hline
\lvert H\rvert&16&20&32&72&144\\
{\lvert X:H\rvert}&45&36&45&35&35.\\
\end{array}
\]

For \(n\ge13\), the factor-free action on 6-subsets is proved in
\cref{prop:alternating-screen}.
\subsection{Classical groups}\label{app:classical}

We use Theorem 2.15 and Appendix Tables A.1--A.7 of Xia--Li
\cite{xiali-solvable}, a restatement of the LPS classification of
nontrivial maximal factorizations with classical socle, together with the
two corrected \(\mathrm P\Omega_8^+(4)\) and
\(\mathrm P\Omega_8^+(16)\) rows of Gill--Giudici--Spiga
\cite[Section 2]{ggs-szep}.  In the table below,
\(V\) is invariant up to \(S\)-conjugacy under \(X\).  Applying
\cref{lem:factor-screen}, inspection of the indicated factorization table
shows that no factor intersection contains \(V\).

We follow the table notation of Xia--Li: \(P_i\) is the standard maximal
parabolic associated with node \(i\), \(P_{i,j}=P_i\cap P_j\) is the
intersection of the indicated standard parabolics containing a common
Borel subgroup, and \(N_i\) denotes the relevant nonsingular \(i\)-space
stabilizer.  The phrase ``full \(2+2\) decomposition stabilizer'' includes
the element interchanging the two summands.

\begin{longtable}{>{\raggedright\arraybackslash}p{.31\textwidth}
                  >{\raggedright\arraybackslash}p{.23\textwidth}
                  >{\raggedright\arraybackslash}p{.34\textwidth}}
\toprule
Socle and conditions & screen subgroup \(V\) & table check \\
\midrule
\endhead
\(\PSL_2(q)\), excluding the allowed and alternating cases
& split-torus normalizer
& not in the nonsplit-torus factor of A.1:1 or in \(P_1\); A.1:5--8 are
checked below \\
\(\PSL_3(q)\), \(q\ne4\), excluding \(q=2\)
& split-torus normalizer
& not in either parabolic or the extension-field factor in A.1:1 \\
\(\PSL_3(4)\), every coordinate group \(X\)
& a pinned maximal subgroup of \(X\)
& complete Table-of-Marks certificate in \cref{app:computer} \\
\(\PSL_n(q)\), even \(n\ge4\)
& \(P_{n/2}\)
& absent from A.1:1--4 (and A.1:5--10 have other dimensions) \\
\(\PSL_n(q)\), odd \(n\ge5\), no graph
& \(P_2\)
& absent from A.1:1--4,10 except \((n,q)=(5,2)\) \\
\(\PSL_n(q)\), odd \(n\ge5\), graph present
& \(P_{2,n-2}\)
& only maximal parabolic overgroups \(P_2,P_{n-2}\); A.1:1--4,10 are
excluded except \((5,2)\) \\
\(\PSp_{2m}(q)\), \(m\ge3\)
& \(P_2\)
& absent from A.2:1--16 (A.2:13 has \(P_2\) only when \(m=2\)) \\
\(\PSp_4(q)\), \(q\) odd
& nonsingular \(2+2\) decomposition stabilizer
& absent from the \(m=2\), odd-\(q\) rows A.2:1,13 \\
\(\PSU_3(q)\)
& \(N_1\)
& absent from the \(n=3\) rows A.3:5--7 \\
\(\PSU_4(q)\), \(q\ge3\)
& full nonsingular \(2+2\) decomposition stabilizer
& absent from A.3:1--4,9; in particular A.3:4 covers \(\PSU_4(4)\) \\
\(\PSU_n(q)\), \(n\ge5\)
& \(N_2\)
& absent from A.3:1--4,10--12 \\
\(\Omega_{2m+1}(q)\), \(m\ge3, q\) odd
& \(P_2\)
& absent from A.4:1--8 \\
\(\mathrm P\Omega^-_{2m}(q)\), \(m\ge4\)
& \(P_2\)
& absent from A.5:1--5 \\
\(\mathrm P\Omega^+_{2m}(q)\), \(m\ge5\)
& \(P_2\)
& absent from A.6:1--13 \\
\(\mathrm P\Omega^+_8(q)\), triality allowed
& central-node \(P_2\)
& fixed by triality and absent from A.7:1--15 \\
\bottomrule
\end{longtable}

The two corrected orthogonal factorizations have socle intersections
\(N_2^-\) and a triality image of
\(\Omega_8^-(q^{1/2})\), for \(q=4,16\).  Neither contains the central-node
maximal parabolic \(P_2\), so they do not alter the last row of the screen.

Here a standard maximal parabolic can only be contained in a proper factor
intersection if it is that intersection.  A subgroup containing
\(P_{2,n-2}\) contains a Borel and is therefore parabolic; its only proper
overgroups are \(P_2\) and \(P_{n-2}\).  These observations justify the
containment, rather than mere equality, assertions in the table.  The
displayed \(N_i\) and full \(2+2\) decomposition stabilizers are standard
maximal geometric subgroups for the stated parameters
\cite[Chapters 3--4]{kleidman-liebeck}; the exceptional
\(\PSU_4(2)\) parameter is routed through the low-rank isomorphism below.

For the first row, the split-torus normalizer has order
\(2(q-1)/(2,q-1)\) and interchanges the two fixed points of its torus, so it
is not contained in \(P_1\).  Nor can it lie in the nonsplit-torus
normalizer: its order would force \(q-1\mid q+1\), apart from nonsimple
parameters.  The other A.1 factors occur only for
\[
q=7,11,16,19,23,29,59.
\]
The values \(7,11\) are allowed.  At \(q=19,29,59,23\), divisibility by the
torus-normalizer order excludes \(A_5\) or \(S_4\).  At \(q=16\), a subgroup
of order 30 lies neither in \(D_{34}\) nor in \(A_5\), the latter because it
would have index 2.  This closes the rank-one family.

For \(S=\PSL_3(q)\), let \(d=(3,q-1)\) and let \(V\) normalize a maximally
split torus.  Its class is invariant under diagonal, field, and graph
automorphisms.  The monomial action shows that \(V\) fixes neither a point
nor a hyperplane, and
\[
 \lvert V\rvert=\frac{6(q-1)^2}{d}.
\]
The extension-field factor in row 1 of Table A.1 has intersection with
\(S\) of order \(3(q^2+q+1)/d\).  For \(q=3\), divisibility excludes
containment; for \(q\ge5\), the former order is larger.  The only remaining
table row is the exceptional \(q=4\) row, and all ten almost-simple groups
with that socle are handled by the finite certificate.  The parameter
\(q=2\) gives the allowed group \(\PSL_3(2)\cong\PSL_2(7)\).

The even four-dimensional symplectic family is the sole infinite row not
closed by the factor screen itself.  It is eliminated uniformly by the
prime-degree subfield action in \cref{prop:sp4-screen}; no finite
calculation is used for that family.
The standard low-rank isomorphisms
\[
\begin{aligned}
\PSL_2(4)&\cong\PSL_2(5)\cong A_5,
&\qquad \PSL_2(9)&\cong A_6,\\
\PSL_3(2)&\cong\PSL_2(7),
&\PSL_4(2)&\cong A_8,\\
\PSU_4(2)&\cong\PSp_4(3).&&
\end{aligned}
\]
and the \(B_2=C_2,D_3=A_3,{}^2D_3={}^2A_3\) identifications route every
overlap to a row already considered.

\subsection{Exceptional groups of Lie type}\label{app:exceptional}

Hering--Liebeck--Saxl classify all such factorizations
\cite[Theorems 1--2]{hls-exceptional}.
The only socles admitting a core-free factorization are
\[
G_2(4),\quad G_2(3^f),\quad F_4(2^f),
\]
together with one extra outer factorization for \(G_2(4).2\).  None of the
listed factors is parabolic.  For every exceptional socle outside these
three families, take \(V\) to be a Sylow subgroup in the defining
characteristic.  Its \(S\)-conjugacy class is invariant under every
automorphism, and the Hering--Liebeck--Saxl theorem and
\cref{lem:factor-screen} make the resulting action factor-free.

It remains to screen the three factorable families.  If no graph
automorphism is present, choose a maximal parabolic.  In the graph-fused
cases choose the graph-stable parabolic intersection \(P_1\cap P_2\) in
\(G_2(3^f)\), or \(P_1\cap P_4\) in \(F_4(2^f)\); each contains a common
Borel.  For the graph-fused \(G_2(3^f)\) branch, the chosen subgroup has
\(p\)-part \(q^6\), whereas the displayed Hering--Liebeck--Saxl factors have
\(p\)-part at most \(q^3\).  This bound is not being asserted for
\(G_2(4)\).  In the \(F_4(2^f)\) branch, the corresponding bounds are at
least \(q^{24}\) and at most \(q^{16}\).  Thus no factor contains the
chosen screen subgroup.  The small isomorphisms \(G_2(2)'\cong\PSU_3(3)\) and
\({}^2G_2(3)'\cong\PSL_2(8)\) have already been covered.

\subsection{Sporadic groups}\label{app:sporadic}

Giudici classifies all factorizations with sporadic socle
\cite[Theorems 1.1--1.3 and Tables 1--4]{giudici}.  Theorem 1.1 and
Tables 1--2 treat simple sporadic groups; Theorem 1.2 and Table 3 give the
new outer cases; and Theorem 1.3 and Table 4 list the exact factorizations.
The factor-free screen uses the complete factorization lists in Tables 1--3,
not the exact-only list in Table 4.
Comparing those loci with the Atlas maximal-subgroup lists
\cite{atlas,atlas-online} gives the following factor-free maximal subgroups.
A row for an outer extension also excludes the factorizations lifted from
the simple socle by Giudici's Lemma 2.1(5).

\begin{center}
\begin{tabular}{ll@{\qquad}ll}
\toprule
\(X\)&\(H\)&\(X\)&\(H\)\\
\midrule
\(M_{11}\)&\(S_5\)&\(M_{12}\)&\(2^3{:}S_4\)\\
\(M_{12}.2\)&\(3^{1+2}{:}D_8\)&\(M_{22}.2\)&\(A_6.2^2\)\\
\(M_{23}\)&\(A_8\)&\(J_2\)&\(3.A_6.2_2\)\\
\(J_2.2\)&\(3.A_6.2^2\)&\(HS\)&\(L_3(4).2_1\)\\
\(HS.2\)&\(L_3(4).2^2\)&\(He\)&\(2^2.L_3(4).S_3\)\\
\(He.2\)&\(2^2.L_3(4).D_{12}\)&\(Ru\)&\(2^6{:}U_3(3){:}2\)\\
\(Suz\)&\(J_2.2\)&\(Suz.2\)&\(J_2.2\times2\)\\
\(Fi_{22}\)&\(O_7(3)\)&\(Fi_{22}.2\)&\(G_2(3).2\)\\
\(Co_1\)&\(2^{11}{:}M_{24}\)&&\\
\bottomrule
\end{tabular}
\end{center}

The simple group \(M_{22}\), and every sporadic almost-simple group not
needed in this display other than \(M_{24}\), has no nontrivial core-free
factorization.  For such a group, apply \cref{lem:factor-screen} to a Sylow
subgroup of the socle;
its conjugacy class is automorphism-invariant and its maximal overgroup has
nontrivial socle intersection.  The rows available in TomLib are separately
cross-checked in \cref{app:computer}.

The sole factor-cover survivor is \(M_{24}\).  Let it act on the 2-subsets of
its natural 24-point set.  The point stabilizer \(M_{22}.2\) is maximal, the
degree is 276, and every involution fixes the support of one of its 2-cycles.
Thus the action is 2-elusive and \cref{lem:elusive-wreath} applies.

For the eight displayed rows not present in the TomLib cross-check, the
maximality assertion is pinned to the ``Maximal subgroups'' section of the
official online ATLAS pages:
\href{https://brauer.maths.qmul.ac.uk/Atlas/v3/spor/He/#maxes}{\(He,He.2\)},
\href{https://brauer.maths.qmul.ac.uk/Atlas/v3/spor/Ru/#maxes}{\(Ru\)},
\href{https://brauer.maths.qmul.ac.uk/Atlas/v3/spor/Suz/#maxes}{\(Suz,Suz.2\)},
\href{https://brauer.maths.qmul.ac.uk/Atlas/v3/spor/F22/#maxes}{\(Fi_{22},Fi_{22}.2\)},
and \href{https://brauer.maths.qmul.ac.uk/Atlas/v3/spor/Co1/#maxes}{\(Co_1\)}.



The four subsections cover every CFSG family, every almost-simple coordinate
group, and every low-rank overlap.  This completes the detailed audit behind
\cref{prop:classification}.

\section{Finite Table-of-Marks certificates}\label{app:computer}
This appendix records the exact machine interface.  The proof uses finite certificates for five small alternating coordinate
groups and for the ten almost-simple coordinate groups with socle
\(\PSL_3(4)\).  Nine sporadic rows are also cross-checked.

For a complete table of marks \(T\), let \(h,c,i\) represent conjugacy
classes of subgroups \(H,C,I\).  The command
\texttt{IntersectionsTom(T,h,c)[i]} is positive exactly when suitable
conjugates of \(H,C\) have intersection in class \(i\).  Such a pair
factors \(X\) precisely when
\[
\lvert H\rvert\,\lvert C\rvert=\lvert X\rvert\,\lvert I\rvert.
\]
A subgroup \(C\) is core-free in an almost-simple \(X\) precisely when its
class does not contain the socle class.  For each pinned named group, the GAP
producer locates the unique least nontrivial normal class from the TomLib
class lengths and orders; the named almost-simple group identifies that class
with its simple socle \(S\).  The producer then checks every class not
containing \(S\) and every intersection class \(i\), verifies
\(\lvert H\cap S\rvert=
  \lvert H\rvert\lvert S\rvert/\lvert X\rvert>1\),
and fails on missing data.  The Python checker validates the pinned inventory,
metadata, arithmetic, and file hash, but does not recompute the TomLib
intersection predicate.

The proof rows are:
\[
\begin{array}{c|ccccc}
X&A_6.2_2&A_6.2_3&A_6.2^2&A_7&S_7\\ \hline
\lvert H\rvert&16&20&32&72&144\\
{\lvert X:H\rvert}&45&36&45&35&35.\\
\end{array}
\]
Since \(\Out(\PSL_3(4))\cong D_{12}\), the following ten TomLib names
exhaust, up to isomorphism, the groups between the socle and its full
automorphism group:
\[
\begin{array}{c|r|r@{\qquad}c|r|r}
X&\lvert H\rvert&\lvert X:H\rvert&X&\lvert H\rvert&\lvert X:H\rvert\\ \hline
L_3(4)&960&21&L_3(4).2_1&384&105\\
L_3(4).2_2&1920&21&L_3(4).2_3&720&56\\
L_3(4).3&216&280&L_3(4).2^2&768&105\\
L_3(4).3.2_2&432&280&L_3(4).3.2_3&1152&105\\
L_3(4).6&1152&105&L_3(4).D_{12}&2304&105.
\end{array}
\]
Every row has zero core-free factor classes.  The producer and integrity
checker are, respectively,
\begin{verbatim}
gap --quitonbreak -q gap/generate-factor-free-scan.g
python3 tests/check-factor-free-scan.py
\end{verbatim}
in the accompanying repository.  The environment is GAP 4.15.1 with
TomLib 1.2.11.  A separate GAP 4.15.1/AtlasRep 2.1.11 regression test
reconstructs the corrected \(q=4\) subfield-normalizer instance from
\cref{prop:sp4-screen}.  All GAP commands
use \texttt{--quitonbreak}; the Python checkers are also run under
\texttt{python3 -O}, so their correctness does not depend on assertions.
The generated TSV has SHA-256
\begin{center}
{\footnotesize
\nolinkurl{9b131720d41ef945a0696794c0493ae9c07e166d4d1f37b054a1d55f0c4837ae}}
\end{center}
and records all twenty-four rows: fifteen proof-critical closeouts and nine
sporadic cross-checks.  This calculation is an exact finite verification
relative to the complete named tables, not an extrapolation from a library
scan.


\section{Dependency, scope, and correction ledger}\label{app:dependencies}

\begingroup\small
\noindent\textbf{CFSG.}
\emph{Exact locus:} the classification theorem.
\emph{Consumed statement:} every nonabelian finite simple socle lies in the
alternating, Lie-type, or sporadic families used in Appendix A.

\smallskip\noindent\textbf{Levchuk--Likharev \cite{levchuk-likharev}.}
\emph{Exact loci:} Theorem 1 and Lemmas 4, 6.
\emph{Consumed statement:} the three listed simple groups have \(\CMP\),
which proves occurrence.

\smallskip\noindent\textbf{Scott \cite{scott}.}
\emph{Exact locus:} the lemma on page 328.
\emph{Consumed statement:} a subdirect subgroup of a direct product of
nonabelian simple groups is a product of disjoint full diagonal strips.

\smallskip\noindent\textbf{Liebeck--Praeger--Saxl
\cite{lps-transitive}\(^{\dagger}\).}
\emph{Attributed locus:} Corollary 3(iv).
\emph{Consumed assumption:} a regular subgroup in the constructed primitive
product action forces a core-free transitive coordinate subgroup.

\smallskip\noindent\textbf{Liebeck--Praeger--Saxl
\cite{lps-maxfact}\(^{\dagger}\).}
\emph{Attributed locus:} Theorem D and Remark 2.
\emph{Consumed assumption:} the large alternating-group factorization screen.

\smallskip\noindent\textbf{Xia--Li \cite{xiali-solvable} and
Gill--Giudici--Spiga \cite{ggs-szep}.}
\emph{Exact loci:} Theorem 2.15 and Tables A.1--A.7; Section 2.
\emph{Consumed statements:} the readable classical-table restatement and
the two corrected orthogonal rows.

\smallskip\noindent\textbf{Hering--Liebeck--Saxl
\cite{hls-exceptional}.}
\emph{Exact locus:} Theorems 1--2.
\emph{Consumed statement:} exceptional Lie-type factorization exhaustiveness.

\smallskip\noindent\textbf{Giudici \cite{giudici}.}
\emph{Exact loci:} Theorem 1.1 and Tables 1--2; Theorem 1.2 and Table 3.
\emph{Consumed statements:} the complete simple-sporadic and new outer
factorization lists used for factor-freeness.  Theorem 1.3 and Table 4 list
only exact factorizations and are not used for that exclusion.

\smallskip\noindent\textbf{Kleidman--Liebeck
\cite{kleidman-liebeck}.}
\emph{Source scope:} Chapters 3--4; the group-specific table routing is
recorded in Appendix~\ref{app:classification}.
\emph{Consumed statement:} maximality of the geometric \(P_i\), \(N_i\),
and decomposition-stabilizer screens used in the classical table.

\smallskip\noindent\textbf{Atlas and online ATLAS
\cite{atlas,atlas-online}.}
\emph{Exact loci:} the maximal-subgroup lists and the five group-page links
after the sporadic table.
\emph{Consumed statement:} maximality of the eight published-only sporadic
screen subgroups.

\smallskip\noindent\textbf{Bray--Holt--Roney-Dougal \cite{bhrd} and
Burness \cite{burness-base}.}
\emph{Exact loci:} Table 8.14; Proposition 4.2 and Table 3.
\emph{Consumed statement:} maximality of the prime-degree
\(\Sp_4(q_0)\) subfield subgroup.

\smallskip\noindent\textbf{Harper \cite{harper} and
Liebeck--Praeger--Saxl \cite{lps-regular}.}
\emph{Exact loci:} Harper's setup preceding Lemma 2.1, Lemma 2.1, and
Theorem 4(a); LPS Lemmas 2.1--2.2 and Corollary 2.3.
\emph{Consumed statements:} the outer automorphisms of
\(\PSp_4(2^f)\) and its three involution classes.

\smallskip\noindent\textbf{GAP/TomLib.}
\emph{Exact environment:} GAP 4.15.1, TomLib 1.2.11; commands in
\cref{app:computer}.
\emph{Consumed output:} fifteen proof-critical factor-free closeouts; the
nine sporadic rows are non-proof cross-checks.

\smallskip\noindent\textbf{Python integrity checker.}
\emph{Exact program:} \texttt{check-factor-free-scan.py}.
\emph{Consumed output:} pinned inventory, producer metadata, hashes, and
arithmetic.  It does not repeat TomLib's negative intersection search.

\smallskip\noindent\textbf{GAP/AtlasRep.}
\emph{Exact environment and program:} GAP 4.15.1, AtlasRep 2.1.11;
\texttt{test-sp4-subfield.g}.
\emph{Consumed output:} none in the theorem; this is a non-proof \(q=4\)
regression for the uniform symplectic construction.
\endgroup

The universal conclusion is conditional on CFSG and the exhaustive published
classification inputs in this ledger.  No property of a section is inferred
from \(\CMP\); only quotient closure is used.  No formal proof assistant is
part of the evidence chain.

The daggered LPS loci are load-bearing source inputs.  Their exact article
texts could not be reopened from an accessible copy during the source traces
through 31 August 2026; the bibliographic locators and the available
audited-copy hash in the repository's reference ledger identify the source
records.
Accordingly, the two consequences
formulated in this ledger are explicit external assumptions at the verification
boundary.  If either original statement has narrower hypotheses, the
corresponding product-action or large-alternating branch remains unchecked;
no such mismatch is presently known.

Versions 1.0.0 and 1.0.1 used an orthogonal-normalizer assertion that fails
for \(X=\Aut(\PSp_4(4))\): the old subgroup has index \(272\), and adjoining
the socle gives a subgroup of index \(2\).  The current proof does not use
that assertion.  It uses the prime-degree subfield subgroup in
\cref{prop:sp4-screen}; at the same parameter the corrected subgroup has
order \(2880\), socle intersection \(720\), and index \(1360\).

The July 2026 Kourovka Notebook prints Problem 18.68 without a solution
annotation.  This dated bibliographic observation is not used in the proof
and does not rule out unindexed or contemporaneous work.

\section*{Generative-AI disclosure}
Generative-AI tools were used for literature search, proof exploration, code
generation, drafting, and critical checking; the author checked and takes
responsibility for the manuscript's arguments, citations, and computations.

\begin{thebibliography}{99}
\small

\bibitem{aschbacher-scott}
M. Aschbacher and L. Scott,
Maximal subgroups of finite groups,
\emph{J. Algebra} \textbf{92} (1985), 44--80,
\url{https://doi.org/10.1016/0021-8693(85)90145-0}.

\bibitem{bhrd}
J.~N. Bray, D.~F. Holt, and C.~M. Roney-Dougal,
\emph{The Maximal Subgroups of the Low-Dimensional Finite Classical Groups},
London Math. Soc. Lecture Note Ser. \textbf{407},
Cambridge University Press, Cambridge, 2013,
\url{https://doi.org/10.1017/CBO9781139192576}.

\bibitem{burness-base}
T.~C. Burness,
On base sizes for actions of finite classical groups,
\emph{J. London Math. Soc.} (2) \textbf{75} (2007), 545--562,
\url{https://doi.org/10.1112/jlms/jdm033}.

\bibitem{atlas}
J.~H. Conway, R.~T. Curtis, S.~P. Norton, R.~A. Parker, and R.~A. Wilson,
\emph{Atlas of Finite Groups}, Clarendon Press, Oxford, 1985.

\bibitem{ggs-szep}
N. Gill, M. Giudici, and P. Spiga,
A generalization of Szep's conjecture for almost simple groups,
\emph{Vietnam J. Math.} \textbf{52} (2024), 325--359,
\url{https://doi.org/10.1007/s10013-023-00635-1}.

\bibitem{giudici}
M. Giudici,
Factorisations of sporadic simple groups,
\emph{J. Algebra} \textbf{304} (2006), 311--323,
\url{https://doi.org/10.1016/j.jalgebra.2006.04.019}.

\bibitem{harper}
S. Harper,
Totally deranged elements of almost simple groups and invariable generating
sets,
\emph{J. London Math. Soc.} (2) \textbf{109} (2024), no.~6, e12935,
\url{https://doi.org/10.1112/jlms.12935}.

\bibitem{hls-exceptional}
C. Hering, M.~W. Liebeck, and J. Saxl,
The factorizations of the finite exceptional groups of Lie type,
\emph{J. Algebra} \textbf{106} (1987), 517--527,
\url{https://doi.org/10.1016/0021-8693(87)90013-5}.

\bibitem{kourovka21}
E.~I. Khukhro and V.~D. Mazurov, eds.,
\emph{The Kourovka Notebook: Unsolved Problems in Group Theory},
21st ed., 2026,
\url{https://kourovkanotebookorg.wordpress.com/wp-content/uploads/2026/07/21tkt.pdf}.

\bibitem{kleidman-liebeck}
P.~B. Kleidman and M.~W. Liebeck,
\emph{The Subgroup Structure of the Finite Classical Groups},
London Math. Soc. Lecture Note Ser. \textbf{129},
Cambridge University Press, Cambridge, 1990,
\url{https://doi.org/10.1017/CBO9780511629235}.

\bibitem{kovacs-composite}
L.~G. Kov\'acs,
Maximal subgroups in composite finite groups,
\emph{J. Algebra} \textbf{99} (1986), 114--131,
\url{https://doi.org/10.1016/0021-8693(86)90058-X}.

\bibitem{kovacs-product}
L.~G. Kov\'acs,
Primitive subgroups of wreath products in product action,
\emph{Proc. London Math. Soc.} (3) \textbf{58} (1989), 306--322,
\url{https://doi.org/10.1112/plms/s3-58.2.306}.

\bibitem{levchuk-likharev}
V.~M. Levchuk and A.~G. Likharev,
Finite simple groups with complemented maximal subgroups,
\emph{Siberian Math. J.} \textbf{47} (2006), 659--668,
\url{https://doi.org/10.1007/s11202-006-0077-7}.

\bibitem{li-wang-xia}
C.~H. Li, L. Wang, and B. Xia,
The exact factorizations of almost simple groups,
\emph{J. London Math. Soc.} (2) \textbf{108} (2023), 1417--1447,
\url{https://doi.org/10.1112/jlms.12784}.

\bibitem{xiali-solvable}
C.~H. Li and B. Xia,
\emph{Factorizations of Almost Simple Groups with a Solvable Factor, and
Cayley Graphs of Solvable Groups},
Mem. Amer. Math. Soc. \textbf{279} (2022), no.~1375,
\url{https://doi.org/10.1090/memo/1375}.

\bibitem{lps-maxfact}
M.~W. Liebeck, C.~E. Praeger, and J. Saxl,
\emph{The Maximal Factorizations of the Finite Simple Groups and Their
Automorphism Groups},
Mem. Amer. Math. Soc. \textbf{86} (1990), no.~432,
\url{https://bookstore.ams.org/memo-86-432}.

\bibitem{lps-transitive}
M.~W. Liebeck, C.~E. Praeger, and J. Saxl,
Transitive subgroups of primitive permutation groups,
\emph{J. Algebra} \textbf{234} (2000), 291--361,
\url{https://doi.org/10.1006/jabr.2000.8547}.

\bibitem{lps-regular}
M.~W. Liebeck, C.~E. Praeger, and J. Saxl,
Regular subgroups of primitive permutation groups,
\emph{Mem. Amer. Math. Soc.} \textbf{203} (2010), no.~952,
\url{https://doi.org/10.1090/S0065-9266-09-00569-9}.

\bibitem{maslova-hall-factors}
N.~V. Maslova,
Nonabelian composition factors of a finite group whose all maximal subgroups
are Hall,
\emph{Siberian Math. J.} \textbf{53} (2012), no.~5, 853--861,
\url{https://doi.org/10.1134/S0037446612050102}.

\bibitem{maslova-revin}
N.~V. Maslova and D.~O. Revin,
Finite groups whose maximal subgroups have the Hall property,
\emph{Siberian Adv. Math.} \textbf{23} (2013), 196--209,
\url{https://doi.org/10.3103/S105513441303005X}.

\bibitem{monakhov-sokhor}
V.~S. Monakhov and I.~L. Sokhor,
On indices of maximal chains in finite groups,
\emph{Results Math.} \textbf{80} (2025), article 155,
\url{https://doi.org/10.1007/s00025-025-02468-5}.

\bibitem{scott}
L.~L. Scott,
Representations in characteristic \(p\),
in \emph{The Santa Cruz Conference on Finite Groups},
Proc. Sympos. Pure Math. \textbf{37}, American Mathematical Society,
Providence, 1980, 319--331,
\url{https://doi.org/10.1090/pspum/037/604599}.

\bibitem{atlas-online}
R.~A. Wilson et al.,
\emph{ATLAS of Finite Group Representations}, version 3,
\url{https://brauer.maths.qmul.ac.uk/Atlas/v3/}.

\end{thebibliography}

\end{document}
